% TODO make hw stylesheet
\documentclass[10pt]{scrartcl}
\usepackage[a4paper, total={6in, 8in}]{geometry}
\usepackage[usenames,dvipsnames,svgnames]{xcolor}
\usepackage[shortlabels]{enumitem}
\usepackage[framemethod=TikZ]{mdframed}
\usepackage{amsmath,amssymb,amsthm}
\usepackage[colorlinks]{hyperref}
\usepackage{multicol}
\usepackage{tabularx}

\hypersetup{
	colorlinks=true,
	linkcolor=blue,
	filecolor=magenta,      
	urlcolor=magenta,
	pdftitle={MAT 314 HW}
}

\usepackage{microtype}
\usepackage{mathtools}
\usepackage[headsepline]{scrlayer-scrpage}
\usepackage{thmtools}
\usepackage[textsize=footnotesize]{todonotes}

\mdfdefinestyle{mdbluebox}{roundcorner=10pt,innerbottommargin=9pt,
	linecolor=blue,backgroundcolor=TealBlue!5,}
\declaretheoremstyle[headfont=\sffamily\bfseries\color{MidnightBlue},
mdframed={style=mdbluebox},]{thmbluebox}

\mdfdefinestyle{mdbloobox}{frametitlefont=\bfseries,innerbottommargin=8pt,
	nobreak=true,backgroundcolor=TealBlue!5,linecolor=blue,}
\declaretheoremstyle[headfont=\bfseries\color{MidnightBlue},
mdframed={style=mdbloobox},headpunct={\\[3pt]},postheadspace=0pt,]{thmbloobox}

\mdfdefinestyle{mdredbox}{frametitlefont=\bfseries,innerbottommargin=8pt,
	nobreak=true,backgroundcolor=Salmon!5,linecolor=RawSienna,}
\declaretheoremstyle[headfont=\bfseries\color{RawSienna},
mdframed={style=mdredbox},headpunct={\\[3pt]},postheadspace=0pt,]{thmredbox}

\mdfdefinestyle{mdgreenbox}{linecolor=ForestGreen,backgroundcolor=ForestGreen!5,
	linewidth=2pt,rightline=false,leftline=true,topline=false,bottomline=false,}
\declaretheoremstyle[headfont=\bfseries\sffamily\color{ForestGreen!70!black},
mdframed={style=mdgreenbox},headpunct={ --- },]{thmgreenbox}

\mdfdefinestyle{mdorangebox}{linecolor=RedOrange,backgroundcolor=RedOrange!5,
	linewidth=2pt,rightline=false,leftline=true,topline=false,bottomline=false,}
\declaretheoremstyle[headfont=\bfseries\sffamily\color{RedOrange!70!black},
mdframed={style=mdorangebox},headpunct={ --- },]{thmorangebox}

\mdfdefinestyle{mdblackbox}{linecolor=black,backgroundcolor=RedViolet!5!gray!5,
	linewidth=3pt,rightline=false,leftline=true,topline=false,bottomline=false,}
\declaretheoremstyle[mdframed={style=mdblackbox}]{thmblackbox}

\declaretheoremstyle[spaceabove=6pt,spacebelow=6pt]{basehead}

\declaretheorem[style=basehead,name=Problem]{problem}
\declaretheorem[style=basehead,name=Problem,numbered=no]{problem*}
\declaretheorem[style=thmbloobox,name=Setup]{setup} % for config geo
\declaretheorem[style=thmredbox,name=Example,sibling=problem]{example}
\declaretheorem[style=thmredbox,name=$\bigstar$ Problem,sibling=problem]{niceprob}
\declaretheorem[style=thmbluebox,name=Theorem,sibling=problem]{theorem}
\declaretheorem[style=thmbluebox,name=Corollary,sibling=theorem]{corollary}
\declaretheorem[style=thmbluebox,name=Proposition,sibling=theorem]{prop}
\declaretheorem[style=thmbluebox,name=Lemma,numberwithin=problem]{lemma}
\declaretheorem[style=thmbluebox,name=Theorem,numbered=no]{theorem*}
\declaretheorem[style=thmbluebox,name=Lemma,numbered=no]{lemma*}
\declaretheorem[style=thmgreenbox,name=Claim,numberwithin=problem]{claim}
\declaretheorem[style=thmgreenbox,name=Claim,numbered=no]{claim*}
\declaretheorem[style=thmblackbox,name=Remark,numberwithin=problem]{remark}
\declaretheorem[style=thmblackbox,name=Remark,numbered=no]{remark*}
\declaretheorem[style=thmgreenbox,name=Definition,sibling=theorem]{definition}
\declaretheorem[style=thmgreenbox,name=Definition,numbered=no]{definition*}
\declaretheorem[style=thmorangebox,name=Warning,sibling=theorem]{warning}
\declaretheorem[style=thmorangebox,name=Warning,numbered=no]{warning*}
\declaretheorem[style=thmorangebox,name=Note,numbered=no]{note}
\declaretheorem[style=thmblackbox,name=Exercise,sibling=problem]{exercise}
\declaretheorem[style=thmblackbox,name=Exercise,numbered=no]{exercise*}
\declaretheorem[style=thmblackbox,name=Question,sibling=problem]{question}
\declaretheorem[style=thmblackbox,name=Question,numbered=no]{question*}

\addtolength{\textheight}{3.14cm}
\providecommand{\ol}{\overline}
\providecommand{\ul}{\underline}
\providecommand{\eps}{\varepsilon}
\providecommand{\half}{\frac{1}{2}}
\providecommand{\dang}{\measuredangle} %% Directed angle
\providecommand{\CC}{\mathbb C}
\providecommand{\FF}{\mathbb F}
\providecommand{\NN}{\mathbb N}
\providecommand{\QQ}{\mathbb Q}
\providecommand{\RR}{\mathbb R}
\providecommand{\ZZ}{\mathbb Z}
\providecommand{\dg}{^\circ}
\providecommand{\ii}{\item}
\providecommand{\setdiff}{\smallsetminus}
\providecommand{\alert}[1]{{\sffamily\textbf{\textcolor{blue}{#1}}}}
\providecommand{\inv}{^{-1}}
\providecommand{\mb}[1]{\mathbf{#1}}
\providecommand{\dif}{\;d}

\reversemarginpar

\newenvironment{soln}{\begin{proof}[Solution]}{\end{proof}}
\newenvironment{walkthrough}{\noindent \textbf{Walkthrough.}}{\medskip}
\newlist{walk}{enumerate}{3}
\setlist[walk]{label=\bfseries (\alph*)}

\providecommand{\pow}{\operatorname{Pow}}
\providecommand{\vocab}[1]{{\textbf{\textcolor{ForestGreen}{#1}}}}
\providecommand{\lcm}{\operatorname{lcm}}
\providecommand{\abs}[1]{\left|#1\right|}

\usepackage{asymptote}
\begin{asydef}
	defaultpen(fontsize(10pt));
	unitsize(1cm);
	size(8cm); // set a reasonable default
	usepackage("amsmath");
	usepackage("amssymb");
	settings.tex="pdflatex";
	settings.outformat="pdf";
	// Replacement for olympiad+cse5 which is not standard
	import geometry;
	// recalibrate fill and filldraw for conics
	void filldraw(picture pic = currentpicture, conic g, pen fillpen=defaultpen, pen drawpen=defaultpen)
	{ filldraw(pic, (path) g, fillpen, drawpen); }
	void fill(picture pic = currentpicture, conic g, pen p=defaultpen)
	{ filldraw(pic, (path) g, p); }
	// some geometry
	pair foot(pair P, pair A, pair B) { return foot(triangle(A,B,P).VC); }
	pair orthocenter(pair A, pair B, pair C) { return orthocentercenter(A,B,C); }
	pair centroid(pair A, pair B, pair C) { return (A+B+C)/3; }
	// cse5 abbreviations
	path CP(pair P, pair A) { return circle(P, abs(A-P)); }
	path CR(pair P, real r) { return circle(P, r); }
	pair IP(path p, path q) { return intersectionpoints(p,q)[0]; }
	pair OP(path p, path q) { return intersectionpoints(p,q)[1]; }
	path Line(pair A, pair B, real a=0.6, real b=a) { return (a*(A-B)+A)--(b*(B-A)+B); }
	// cse5 more useful functions
	picture CC() {
		picture p=rotate(0)*currentpicture;
		currentpicture.erase();
		return p;
	}
	pair MP(Label s, pair A, pair B = plain.S, pen p = defaultpen) {
		Label L = s;
		L.s = "$"+s.s+"$";
		label(L, A, B, p);
		return A;
	}
	pair Drawing(Label s = "", pair A, pair B = plain.S, pen p = defaultpen) {
		dot(MP(s, A, B, p), p);
		return A;
	}
	path Drawing(path g, pen p = defaultpen, arrowbar ar = None) {
		draw(g, p, ar);
		return g;
	}
\end{asydef}

\usepackage{mathrsfs}

\ihead{\footnotesize MAT 314 HW03}
\ohead{\footnotesize Last compiled \today}

\title{\vspace{-2em}MAT 314 HW03}
\author{\large Aditya Pahuja}
\date{\large Due 03/06}
\begin{document}
\maketitle

\begin{problem*}[1.2]
	Prove that, for $n\ne 0$, $\cos(2\pi/n)$ is an algebraic number.
\end{problem*}

\begin{soln}
	We know that
	\[(\cos(2\pi/n) + i\sin(2\pi/n))^n = \cos(2\pi) + i\sin(2\pi) = 1.\]
	Let $x = \cos(2\pi/n)$. Then, the left-hand side expands to
	\[\sum_{k=0}^n \binom nk x^{n-k} (\sqrt{1-x^2})^{k}\cdot i^k.\]
	The real part of this quantity is the sum of the addends
	for which $i^k$ is real, i.e. $k$ is even, which amounts to
	\[\sum_{0\le k'\le n/2} \binom n{2k'}
	x^{n-2k'}(1-x^2)^{k'}\cdot (-1)^{k'} = 1.\]
	Thus $\cos(2\pi/n)$ is a root of an integer polynomial
	\[P(x) = \sum_{0\le k'\le n/2} \binom n{2k'}
	x^{n-2k'}(1-x^2)^{k'}\cdot (-1)^{k'} - 1,\]
	which is what it means for $\cos(2\pi/n)$ to be algebraic.
\end{soln}

\begin{problem*}[2.1]
	For which positive integers $n$ does
	$x^2 + x + 1$ divide $x^4 + 3x^3 + x^2 + 7x + 5$ in $[\ZZ/(n)][x]$?
\end{problem*}

\begin{soln}
	Let $P(x) = x^4 + 3x^3 + x^2 + 7x + 5$. The division algorithm gives
	\[P(x) = (x^2 + x + 1)(x^2 + 2x - 2) + 7x + 7.\]
	We want the remainder to be zero in $[\ZZ/(n)][x]$,
	which means that $7x + 7 = nQ(x)$ for some $Q(x)\in[\ZZ/(n)][x]$.
	The only choice for $n$ here is thus $\boxed 7$.
\end{soln}

\begin{problem*}[3.5]
	The derivative of a polynomial $f$ with coefficients in a field $F$
	is defined by the formula
	\[(a_nx^n + \cdots + a_1x + a_0)' = na_nx^{n-1} + \cdots + 1a_1.\]
	\begin{enumerate}[(a)]
		\ii Prove the product rule $(fg)' = f'g + fg'$ and the chain rule
		$(f\circ g)' = (f'\circ g)\cdot g'$.
		\ii Let $\alpha$ be an element of $F$. Prove that $\alpha$
		is a multiple root of a polynomial $f$ if and only if
		it is a common root of $f$ and $f'$.
	\end{enumerate}
\end{problem*}

\begin{soln}
	(a) Let $f(x) = a_nx^n + \cdots + a_1x+a_0$ and
	$g(x) = b_mx^m + \cdots + b_1x + b_0$.
	Then we observe that for any monomial $cx^k$,
	\begin{align*}
		(f(x)\cdot cx^k)' &=
		(ca_nx^{k+n} + \cdots + ca_1x^{k+1} + ca_0x^k)'\\
		&= \sum_{j=0}^n c(k+j)a_j x^{k+j-1}\\
		&= cx^k\sum_{j=0}^n(ja_j x^{j-1}) +
		ckx^{k-1} \sum_{j=0}^n a_jx^j\\
		&= (cx^k)f'(x) + (cx^k)' f(x),
	\end{align*}
	so the product rule holds when one of the polynomials is a monomial.
	We can clearly see that the derivative is a linear operator on $F[x]$,
	so this gives us
	\begin{align*}
		(fg)' = \sum_{j=0}^m (f(x)\cdot b_jx^j)'
		&= \sum_{j=0}^m f'(x)b_jx^j + f(x)\cdot jb_jx^{j-1}\\
		&= f'(x)\sum_{j=0}^m b_jx^j + f(x)\sum_{j=0}^m jb_jx^{j-1}\\
		&= f'g + fg'.
	\end{align*}

	We now prove the chain rule. First, we will prove the chain rule
	when $f(x) = cx^k$ is a monomial by induction on $k$.
	If $k = 0$ then $f$ is constant, so $f\circ g$ is constant
	and hence its derivative is zero, which aligns with $f'\circ g = 0$.
	Now suppose that the chain rule holds for monomials with degree $k-1$.
	Then if $f(x) = cx^k$,
	\begin{align*}
		(f\circ g)' = (c\cdot g^k)'
		&= c(g^{k-1}\cdot g)'\\
		&= c\cdot ((g^{k-1})'\cdot g + g^{k-1}\cdot g')\\
		&= c\cdot ((k-1)g^{k-2}\cdot g'\cdot g + g^{k-1}\cdot g')\\
		&= c\cdot kg^{k-1}\cdot g',
	\end{align*}
	where the second-to-last equality comes from applying
	the chain rule to the composition $x\mapsto g(x)\mapsto g(x)^{k-1}$
	via the induction hypothesis. To finish, we can extend to arbitrary $f$
	via linearity:
	\begin{align*}
		(f\circ g)' &= \left(\sum_{j=0}^n a_n g^n\right)'
		= \sum_{j=0}^n (a_n g^n)'\\
		&= \sum_{j=0}^n na_n g^{n-1}\cdot g' = g'\cdot (f'\circ g.)
	\end{align*}

	(b) Suppose $\alpha$ is a multiple root of $f$.
	Then $(x-\alpha)^2\mid f$, so, writing $f(x) = (x-\alpha)^2 g(x)$,
	\[f'(x) = (x-\alpha)^2 g'(x) + 2(x-\alpha)g(x),\]
	which is clearly divisible by $x-\alpha$, hence $f'$ has $\alpha$ as a root.
	
	Conversely, suppose $\alpha$ is a common root of $f$ and $f'$.
	Writing $f(x) = (x-\alpha)g(x)$, we get
	\[f'(x) = (x-\alpha)g'(x) + g(x),\]
	so setting $x=\alpha$ yields
	\[0 = f'(\alpha) = g(\alpha),\]
	hence $x-\alpha\mid g(x)$ implying $(x-\alpha)^2\mid f(x)$ as desired.
\end{soln}

\begin{problem*}[5.7]
	Let $F$ be a field and let $R = F[t]$ be the polynomial ring.
	Let $R'$ be the ring extension $R[x]/(tx-1)$ obtained by
	adjoining an inverse of $t$ to $R$. Prove that this ring can be identified
	as the ring of \emph{Laurent polynomials}, which are
	finite linear combinations of powers of $t$, negative exponents included.
\end{problem*}

\begin{soln}
	Let $F[t,t\inv]$ be the ring of Laurent polynomials in $t$ over $F$.
	We wish to show that $F[t,x]/(tx-1)\cong F[t,t\inv]$,
	which we will do by showing that the map
	$\phi\colon F[t,x]\to F[t,t\inv]$ given by
	\[\phi\left(\sum_{k=0}^n \sum_{i=0}^k a_{i,k-i}t^ix^{k-i}\right)
	= \sum_{k=0}^n \sum_{i=0}^k a_{i,k-i}t^{2i-k}\]
	is a homomorphism of rings with kernel $(tx-1)$.
	We can see easily that $\phi(1) = 1$ and $\phi(p+q) = \phi(p) + \phi(q)$
	for any $p,q\in F[t,x]$. To show that $\phi(pq) = \phi(p)\phi(q)$,
	we only need to prove this in the case that
	one of $p$, $q$ is a monomial (WLOG $q$),
	since we can get at general polynomials by linearity.
	The case where $q$ is a monomial is easy:
	if $q(t,x) = ct^rx^s$, then
	\[\phi(pq) = \phi\left(\sum_{k=0}^n \sum_{i=0}^k
	ca_{i,k-i}t^{i+r}x^{k-i+s}\right)
	= \sum_{k=0}^n \sum_{i=0}^k a_{i,k-i}t^{2i-k}\cdot ct^{r-s}
	= \phi(p)\cdot \phi(q)\]
	since $\phi(q) = ct^{r-s}$. To be more precise about general $q$,
	we note that $\phi(pq)$ can be expanded as a sum of terms
	of the form $\phi(pq_{i,j})$ where $q_{i,j}$ is the monomial
	$b_{i,j}t^ix^j$ in $q$, and then we can write this as
	$\phi(p)\phi(q_{i,j})$ and use the distributive property
	and the additivity of $\phi$ to recombine the $\phi(q_{i,j})$
	into $\phi(q)$.
	
	Now we show that $\ker\phi = (tx-1)$. If
	\[\phi\left(\sum_{k=0}^n \sum_{i=0}^k a_{i,k-i}t^ix^{k-i}\right)
	= \sum_{k=0}^n \sum_{i=0}^k a_{i,k-i}t^{2i-k} = 0,\]
	we can rewrite
	\[\sum_{k=0}^n \sum_{i=0}^k a_{i,k-i}t^ix^{k-i}
	= \sum_{d=1}^n \sum_{i = 0}^{ n-d} a_{d+i,i} t^{d+i}x^i
	+ \sum_{d=1}^n \sum_{i = 0}^{n-d} a_{i,d+i} t^{i}x^{d+i}
	+ \sum_{i=0}^{n} a_{i,i} t^id^i.\]
	Then the image of this polynomial under $\phi$ being zero
	is equivalent to
	\begin{align*}
		\phi\left(\sum_{i = 0}^{ n-d} a_{d+i,i} t^{d+i}x^i\right)
		&= t^d \sum_{i = 0}^{ n-d} a_{d+i,i} = 0\\
		\phi\left(\sum_{i = 0}^{n-d} a_{i,d+i} t^{i}x^{d+i}\right)
		&= x^d \sum_{i = 0}^{ n-d} a_{i,d+i} = 0\\
		\phi\left(\sum_{i=0}^{n} a_{i,i} t^id^i\right)
		&= \sum_{i=0}^{n} a_{i,i} = 0.
	\end{align*}
	I claim that this implies each of the arguments of $\phi$ above
	are divisible by $tx-1$, which is enough to imply that
	the kernel is $(tx-1)$ (since clearly any polynomial with $tx-1$
	as a factor gets sent to zero because $\phi(tx-1) = t\cdot t\inv-1 = 0$).
	Indeed
	\[\sum_{i = 0}^{ n-d} a_{d+i,i} t^{d+i}x^i
	= t^d\cdot \sum_{i = 0}^{ n-d} a_{d+i,i} t^{i}x^i\]
	is $t^d$ times a polynomial in $tx$ whose coefficients sum to zero;
	the theory of single-variable polynomials thus tells us that
	$tx = 1$ being a root implies that $tx - 1$ is a factor of this polynomial.
	Similar logic applies to the other two equations
	(the second after factoring $x^d$ instead),
	so our original polynomial that got sent to zero
	can be written as a sum of polynomials divisible by $tx-1$,
	which implies the result.
\end{soln}

\begin{problem*}[6.8]
	Let $I$ and $J$ be ideals of a ring $R$ such that $I+J = R$.
	\begin{enumerate}[(a)]
		\ii Prove that $IJ = I\cap J$.
		\ii Prove the \emph{Chinese Remainder Theorem}:
		For any pair $a$, $b$ of elements of $R$, there is an element $x$
		such that $x\equiv a\pmod I$ and $x\equiv b\pmod J$.
		\ii Prove that if $IJ = 0$, then $R$ is isomorphic to
		the product ring $(R/I)\times (R/J)$.
		\ii Describe the idempotents corresponding to the
		product decomposition in (c).
	\end{enumerate}
\end{problem*}

\begin{soln}
	(a) On one hand $IJ\subseteq I\cap J$ because if $i\in I$ and $j\in J$
	then $ij\in I$ and $ij\in J$ both hold by definition of ideal.
	On the other hand, suppose that $a\in I\cap J$.
	Then we can write $1 = i + j$ for some $i\in I$, $j\in J$,
	so that $ai + aj\in IJ$ because $ai$ and $aj$ are both
	products of an element of $I$ with an element of $J$.
	
	(b) As before write $1 = i+j$ and set $x = aj + bi$. Then
	\[x\equiv aj = a(1 - i) \equiv a\pmod I,\quad
	x\equiv bi = b(1 - j)\equiv b\pmod J\]
	as desired.
	
	(c) Let $\phi_I\colon R\to R/I$ and $\phi_J\colon R\to R/J$
	be the projection homomophisms.
	Consider the homomorphism $\phi\colon R\to (R/I)\times (R/J)$ given by
	\[\phi(x) := (\phi_I(x), \phi_J(x)).\]
	We have $x\in\ker\phi$ if and only if $\phi_I(x) = \phi_J(x) = 0$,
	so the kernel is $I\cap J = IJ = 0$.
	
	(d) Let $e\in R$ be the element such that $e\equiv 1\pmod I$ and
	$e\equiv 0\pmod J$. We can see that $e$ is idempotent.
	Moreover, $ex \equiv 0\pmod I$ if and only if $x\in I$;
	the only if direction is because if $ex \in I$ and $x = i+j$ for
	$i\in I$, $j\in J$, then $ex\equiv ej\equiv j\pmod I$,
	implying $j\in I\cap J = IJ = 0$ so $x\in I$.
	Thus the kernel of the map $f\colon R\to R/I$ given by $f(x) = ex$
	is $I$, so $R/I\cong eR$ and similarly $R/J\cong (1-e)R$, giving us
	\[R\cong eR\times (1-e)R.\qedhere\]
\end{soln}
















\end{document}